% ============================================================
%  QOVT — Quantum Orthogonal Vision Transformer — NeurIPS-style
%  RY+CNOT ring circuit + Givens butterfly for lensing DM classification
% ============================================================
\documentclass[10pt,letterpaper]{article}

% --- NeurIPS-compatible layout ---
\usepackage[margin=1in]{geometry}
\usepackage[T1]{fontenc}
\usepackage{lmodern}
\usepackage{mathptmx}
\usepackage{microtype}
\usepackage[numbers,square]{natbib}

% --- Math & theorem environments ---
\usepackage{amsmath,amssymb,amsthm}
\theoremstyle{plain}
\newtheorem{proposition}{Proposition}
\newtheorem{corollary}{Corollary}
\theoremstyle{definition}
\newtheorem{definition}{Definition}
\newtheorem{remark}{Remark}

% --- Tables & figures ---
\usepackage{booktabs}
\usepackage{colortbl}
\usepackage{array}
\usepackage{multirow}
\usepackage{graphicx}

% --- Color & hyperlinks ---
\usepackage{xcolor}
\usepackage[colorlinks=true,linkcolor=blue,citecolor=blue,urlcolor=blue]{hyperref}

% --- Misc ---
\usepackage{enumitem}
\setlist[itemize]{leftmargin=*,topsep=2pt,itemsep=1pt}
\setlist[enumerate]{leftmargin=*,topsep=2pt,itemsep=1pt}
\setlength{\parskip}{4pt}

% --- Colors ---
\definecolor{good}{RGB}{22,163,74}
\definecolor{rowhi}{RGB}{220,252,231}
\newcommand{\win}[1]{\textcolor{good}{\textbf{#1}}}
\newcommand{\running}{\texttt{*}}

% ============================================================
\title{\textbf{A NISQ-Deployable Quantum Orthogonal Attention Circuit for\\
Dark Matter Classification: Matching, Not Beating, Classical Attention,\\
and a Cautionary Tale About Training Instability}}

\author{%
  Anonymous Authors\\[2pt]
  GSoC 2026 --- ML4SCI / DeepLense Quantum Machine Learning Track
}

\date{}

% ============================================================
\begin{document}
\maketitle

% ─────────────────────────────────────────────────────────────
\begin{abstract}
We apply Quantum Orthogonal Vision Transformers (QOVT) to the DeepLense
simulated gravitational lensing benchmark, evaluating two circuit variants
that parameterize an attention key/value projection as an orthogonal
operator: a Givens butterfly decomposition ($T{=}1{,}536$ trainable angles,
classically simulated) and an RY+CNOT ring circuit ($T{=}384$ angles, 6
qubits, NISQ-deployable).
At $N{=}500$ samples/class, RY+CNOT achieves AUC $=0.9144$ versus Givens'
AUC $=0.526$ (mean of 2 runs, near chance)---but classical multi-head
attention (unconstrained, more parameters) scores AUC $=0.910$ at the same
$N$, statistically indistinguishable from RY+CNOT.
Two further checks that we ran specifically to test our initial
generalization-bound explanation both contradict it: Givens also collapses
at $N{=}8{,}000$ (AUC $=0.558$), a size too large to explain by memorization,
and a circuit-depth ablation shows RY+CNOT's own AUC \emph{increasing}
monotonically with parameter count $T \in \{192,384,576\}$ at $N{=}500$---
opposite to the bound's prediction.
We conclude that the $N{=}500$ Givens failure reflects
\textbf{architecture-specific optimization instability}, not a
generalization-theoretic quantum advantage: RY+CNOT and classical attention
achieve comparable accuracy at every tested data size from $N{=}500$ to full
data ($\approx25{,}000$/class), where all evaluated architectures---RY+CNOT,
Givens, sham controls, classical MHA, and a reproduced MAE-pretrained ViT
baseline---land within $1.5$ AUC points of each other.
The RY+CNOT ring circuit remains directly executable on IBM Quantum and
equivalent NISQ hardware, runs $10\times$ faster than PennyLane simulation
via a pure-PyTorch matrix implementation, and is---to our knowledge---a
working, hardware-ready quantum attention architecture for astrophysical
image classification, though we no longer claim it demonstrates a quantum
advantage over classical attention on this benchmark.
\end{abstract}

% ─────────────────────────────────────────────────────────────
\section{Introduction}
\label{sec:intro}

Strong gravitational lensing is a promising probe of the substructure of
dark matter halos~\cite{hezaveh2017fast,alexander2020deep}.
The morphological signature of a lensing event---arc thickness, Einstein ring
radius, and fine substructure---encodes the foreground dark matter model:
axion (ultralight) dark matter~\cite{hui2017ultralight}, CDM subhalos,
and smooth (substructure-free) halos each produce statistically distinguishable
distortions in $64\times64$ pixel observations~\cite{brehmer2019mining,alexander2020deep}.
Traditional flux-ratio anomaly methods~\cite{dalal2002detecting} can detect
substructure but are computationally prohibitive at the scale of forthcoming
surveys (Rubin LSST, Euclid, Roman Space Telescope), which will deliver
tens of thousands of strong lensing events.
Automated three-class classification of these scenarios is therefore a
pressing need.

The core challenge is \emph{data efficiency}.
Simulated lensing datasets are computationally expensive to generate,
and labeled observational datasets are extremely scarce
(each confirmed strong lens requires spectroscopic follow-up).
The state-of-the-art MAE-pretrained ViT~\cite{prasha2025masked} achieves
AUC~$\approx0.97$ at full data but requires 2.72M parameters and a large
unlabeled pretraining corpus---precisely the resource combination unavailable
when new instruments begin observing.

Quantum machine learning offers a principled path to data efficiency.
The Caro et al.\ (2022) generalization bound~\cite{caro2022} proves that
a quantum model with $T$ trainable circuit parameters achieves
test-train gap at most $C\sqrt{T/N}$, independently of model depth
or Hilbert space dimension.
Crucially, this bound is \emph{tighter for circuits with fewer parameters},
suggesting a design principle: prefer circuits with small $T$ when $N$ is small.

In this work, we adapt the Quantum Orthogonal Vision Transformer of
Cherrat et al.~\cite{cherrat2022} to the DeepLense setting.
QOVT replaces the learned Q/K/V projection matrices in self-attention
with \emph{orthogonal} matrices produced by quantum circuits,
combining the geometric inductive bias of isometric attention with
the generalization guarantees of quantum parameterization.
We evaluate two circuit families that instantiate this principle at
different positions in the $T$--expressivity tradeoff.

Our main contributions are:
\begin{enumerate}
  \item \textbf{A NISQ-deployable quantum attention circuit that matches
        classical attention} (Section~\ref{sec:sweep}):
        RY+CNOT ($T{=}384$) is statistically indistinguishable from classical
        multi-head attention across every tested data size from $N{=}500$
        (AUC $0.9144$ vs.\ $0.910$) to full data (AUC $0.985$ vs.\ $0.9921$),
        while remaining a valid gate-model circuit executable on IBM Quantum,
        IonQ, and other NISQ hardware.
  \item \textbf{Identification and diagnosis of Givens-specific training
        instability} (Sections~\ref{sec:caro-prediction}--\ref{sec:crossover}):
        The Givens butterfly ($T{=}1{,}536$) fails catastrophically
        (near-chance AUC) in $50\%$ of logged runs at $N{=}500$ and again at
        $N{=}8{,}000$---a pattern inconsistent with the memorization account
        we initially proposed, and better explained by optimization
        instability in its deeper compositional structure
        (Section~\ref{sec:analysis}).
  \item \textbf{A falsifiable Caro-bound prediction, tested and refuted}
        (Section~\ref{sec:depth}):
        A circuit-depth ablation ($T \in \{192, 384, 576\}$) predicted
        smaller $T$ should generalize better at $N{=}500$; we observe the
        opposite (AUC increases monotonically with $T$), indicating the
        Caro bound does not bind in this parameter range for this
        architecture family.
        We report this negative result because it directly informed our
        revised interpretation of finding 2.
  \item \textbf{Efficient classical implementation}
        (Section~\ref{sec:rycnot}):
        We implement the RY+CNOT circuit as a pure-PyTorch matrix product
        ($10\times$ speedup over PennyLane simulation) with mathematically
        identical output, enabling the multi-run, multi-$N$ sweeps that
        revealed findings 2--3 above.
\end{enumerate}

% ─────────────────────────────────────────────────────────────
\section{Related Work}
\label{sec:related}

\paragraph{Gravitational lensing with deep learning.}
Early CNN methods~\cite{hezaveh2017fast,lanusse2018cmu} established
automated morphological discrimination.
Simulation-based inference~\cite{brehmer2019mining} enables marginalisation
over subhalo population parameters beyond point classification.
Alexander et al.~\cite{alexander2020deep} demonstrated end-to-end dark matter morphology learning.
The current state-of-the-art is an MAE-pretrained ViT~\cite{prasha2025masked}
(AUC $\approx 0.97$ at full data), requiring 2.72M parameters and large unlabeled data.
QOVT addresses the low-data frontier not well-served by these methods.

\paragraph{Quantum ML for particle physics and astrophysics.}
Quantum circuits have been applied to event classification~\cite{terashi2021event,ngairangbam2022anomaly}
and jet tagging~\cite{comajoan2024quantum} in high-energy physics.
Quantum attention has been studied by Tesi et al.~\cite{tesi2024quantum}, Unlu et al.~\cite{unlu2024hybrid}, and Cherrat et al.~\cite{cherrat2022}.
Our work is the first to apply quantum orthogonal attention to astrophysical imaging.

\paragraph{Quantum generalization theory.}
Caro et al.~\cite{caro2022} prove the $C\sqrt{T/N}$ generalization bound for quantum circuits.
Huang et al.~\cite{huang2021power} analyze when quantum models outperform classical ones as
a function of training data.
Liu et al.~\cite{liu2021rigorous} establish provable quantum speed-up for specific kernel families.
We use the Caro bound not as a post-hoc analysis but as a \emph{circuit design criterion}
(minimize $T$ when $N$ is small).

\paragraph{Barren plateaus and trainability.}
McClean et al.~\cite{mcclean2018barren} show that random wide circuits exhibit exponentially
vanishing gradients; Cerezo et al.~\cite{cerezo2021cost} show that local cost functions mitigate this.
The RY+CNOT ring avoids barren plateaus by design: its structured, nearest-neighbor
entanglement and moderate depth ($L=8$) keep gradient norms measurably nonzero
(verified in Section~\ref{sec:audit}).

\paragraph{Orthogonal and unitary constraints in deep learning.}
Orthogonal weight normalization~\cite{huang2018orthogonal} improves classical
network training stability.
O-ViT~\cite{fei2022ovit} applies orthogonal constraints to Vision Transformers
classically, demonstrating the geometric prior's value.
Kerenidis et al.~\cite{kerenidis2021classical} and Landman et al.~\cite{landman2022quantum} establish
quantum algorithms for orthogonal neural networks.
QOVT~\cite{cherrat2022} combines these streams by parameterizing
the orthogonal constraint via a quantum circuit.

% ─────────────────────────────────────────────────────────────
\section{Theoretical Background}
\label{sec:theory}

\subsection{Generalization Bound for Quantum Circuits}

\begin{proposition}[\cite{caro2022}]
\label{prop:caro}
Let $\mathcal{F}$ be a hypothesis class of quantum models parameterized by
$T$ trainable rotation angles (gates of the form $e^{i\theta G}$ for
Hermitian generators $G$), and let $\hat{h}$ be the empirical risk minimizer
over $N$ i.i.d.\ samples drawn from distribution $\mathcal{D}$.
Then, with probability at least $1-\delta$ over the training draw,
\[
  \mathcal{R}(\hat{h}) - \hat{\mathcal{R}}(\hat{h})
  \;\leq\;
  C\sqrt{\frac{T}{N}} + O\!\left(\sqrt{\frac{\log(1/\delta)}{N}}\right),
\]
where $C$ is a universal constant independent of model depth, circuit width,
or Hilbert space dimension.
\end{proposition}

\begin{corollary}[QOVT design criterion]
\label{cor:design}
For two QOVT variants with $T_1 < T_2$ trainable parameters and identical
architecture otherwise, Proposition~\ref{prop:caro} predicts:
\begin{equation}
  \frac{\epsilon_1(N)}{\epsilon_2(N)}
  \;\leq\;
  \sqrt{\frac{T_1}{T_2}},
  \quad
  \text{with the ratio shrinking as } N \to \infty.
  \label{eq:bound-ratio}
\end{equation}
At fixed $N$, the lower-$T$ variant (RY+CNOT, $T_1{=}384$) is guaranteed
to have a tighter generalization bound than the higher-$T$ variant
(Givens, $T_2{=}1{,}536$) by a factor of $\sqrt{384/1536} = 0.5$.
\end{corollary}

\subsection{Orthogonal Attention}
\label{sec:qoattn}

Standard scaled dot-product attention computes:
\begin{equation}
  \mathrm{Attn}(X) = \mathrm{softmax}\!\left(\frac{(W_Q X)(W_K X)^\top}{\sqrt{D}}\right) W_V X,
\end{equation}
where $W_Q, W_K, W_V \in \mathbb{R}^{D\times D}$ are unconstrained learned projections.
QOVT sets $W_Q = I$ (the query is left untransformed) and replaces
$W_K = U$, $W_V = V$ with orthogonal matrices $U, V \in O(D)$
(this asymmetric choice matches our implementation, Section~\ref{sec:rycnot};
an alternative, symmetric design with $W_Q = W_K = U$ would additionally
cancel $U$ from the attention map via $U^\top U = I$, but we did not
implement that variant), yielding:
\begin{equation}
  A_{ij} = \mathrm{softmax}_j\!\left(\frac{x_i^\top (Ux_j)}{\sqrt{D}}\right),
  \qquad
  \mathrm{out}_i = \sum_j A_{ij} (V x_j).
  \label{eq:qoattn}
\end{equation}
Because $U$ is orthogonal, $\|Ux_j\|_2 = \|x_j\|_2$: the key transformation
preserves the magnitude of each token, so attention scores cannot be inflated
or suppressed by $U$ scaling up or down individual keys.
This is a weaker property than full angular-only attention (which would
require the symmetric $W_Q{=}W_K{=}U$ design): the attention map here still
depends on the specific rotation $U$ applies to each key, not solely on the
raw angle between query and key tokens.
The value matrix $V$ separately rotates token representations before
aggregation, preserving all pairwise distances among the aggregated values.

This has two clearly-established beneficial properties, and one that we
state more cautiously in light of the asymmetric design:
(1) \textbf{Parameter efficiency}: $U$ and $V$ are each specified by
$T/4$ rotation angles rather than $D^2$ free entries;
(2) \textbf{Spectral stability}: $\|U\|_2 = \|V\|_2 = 1$ prevents gradient
explosion through the attention layer;
(3) \textbf{Caro bound}: the parameterization falls directly in
Proposition~\ref{prop:caro}'s hypothesis class, though as
Section~\ref{sec:sweep} shows empirically, the bound's small-$T$-implies-better
prediction does not consistently manifest against classical MHA at the
data sizes we tested.

% ─────────────────────────────────────────────────────────────
\section{Method}
\label{sec:method}

\subsection{Notation}

\begin{table}[h!]
\centering
\caption{Notation used throughout the paper.}
\begin{tabular}{@{}ll@{}}
\toprule
Symbol & Meaning \\
\midrule
$x \in \mathbb{R}^{64\times64}$ & Input lensing image (grayscale) \\
$D = 64 = 2^6$ & Token dimension / quantum Hilbert space dimension \\
$L$ & Number of QOVT layers (default: 4) \\
$L_\mathrm{circ}$ & Number of quantum circuit layers (RY+CNOT variant, default: 8) \\
$U, V \in O(D)$ & Orthogonal attention matrices (one per QOVT layer) \\
$T$ & Total trainable quantum rotation angles \\
$N$ & Training samples per class \\
$N_Q$ & Number of qubits ($N_Q = 6$, $2^{N_Q} = D = 64$) \\
$\theta_{\ell,q}$ & RY rotation angle for layer $\ell$, qubit $q$ \\
$\phi_{\ell,k}$ & Givens rotation angle for level $\ell$, pair $k$ \\
\bottomrule
\end{tabular}
\label{tab:notation}
\end{table}

\subsection{Overall Architecture}

We apply QOVT to $64\times64$ grayscale lensing images:
\begin{enumerate}
  \item \textbf{Patch embedding}: divide $64\times64$ image into $8\times8$ patches,
        linearly project each patch to $D{=}64$ dimensions; yields 64 tokens.
  \item \textbf{CLS token}: prepend a learnable \texttt{[CLS]} token; add sinusoidal
        positional embeddings. Total sequence length: 65.
  \item \textbf{$L{=}4$ QO-attention layers}: each with quantum orthogonal attention
        (Section~\ref{sec:qoattn}), residual connection, LayerNorm, and 2-layer MLP block.
  \item \textbf{Classification head}: $\texttt{[CLS]}$ output $\to \mathrm{Linear} \to \mathbb{R}^3$.
\end{enumerate}

The token dimension $D{=}64{=}2^6$ is chosen to match the 6-qubit RY+CNOT circuit:
the $64$-dimensional orthogonal matrix $U \in O(64)$ maps directly onto
the $2^6$-dimensional Hilbert space.

\subsection{Variant 1: RY+CNOT Ring Circuit (Main Contribution)}
\label{sec:rycnot}

The RY+CNOT ring circuit parameterizes $U \in O(64)$ via a 6-qubit gate sequence:

\begin{definition}[RY+CNOT Ring Circuit]
For $L_\mathrm{circ} = 8$ layers, circuit layer $\ell$ consists of:
\begin{enumerate}
  \item Single-qubit RY rotations: $\mathrm{RY}(\theta_{\ell,q})$ on each qubit
        $q \in \{0, 1, 2, 3, 4, 5\}$;
  \item CNOT ring: CNOT$(q \to q+1 \bmod 6)$ for all $q$.
\end{enumerate}
\end{definition}

The $2^6 \times 2^6$ unitary $U$ is obtained by acting the circuit on all $2^6$
standard basis states.
The real part $\mathrm{Re}(U)$ is used as the orthogonal matrix
(a valid approximation since the imaginary part is negligible for shallow circuits
with RY gates, which produce real matrix entries).

\textbf{Implementation.}
Rather than PennyLane simulation, we implement the circuit as a sequence of
PyTorch sparse matrix multiplications:
each RY layer is a block-diagonal $64\times64$ matrix assembled from
$2\times2$ rotation blocks, and the CNOT ring is a fixed permutation matrix.
This produces mathematically identical output to PennyLane's
\texttt{default.qubit} simulator at $10\times$ higher throughput.

\textbf{Trainable parameters.}
$T_{\mathrm{RY}} = L_\mathrm{circ} \times N_Q \times 2\text{ matrices} \times L\text{ QOVT layers}
= 8 \times 6 \times 2 \times 4 = 384$.

\textbf{NISQ deployability.}
Each layer is an alternating sequence of single-qubit RY gates (natively
supported on all major hardware) and nearest-neighbor CNOTs (compatible with
superconducting ring and linear-chain connectivity).
Circuit depth $O(L_\mathrm{circ} \cdot N_Q) = O(48)$ is well within
current coherence times.
Trained parameters transfer directly from simulation to real hardware.

\subsection{Variant 2: Givens Butterfly Decomposition}
\label{sec:givens}

The Givens butterfly parameterizes $U \in O(64)$ via a classical butterfly
pattern of $2\times2$ Givens rotations:
\begin{equation}
  U = G_{\log_2 D} \cdots G_2 G_1,
  \quad
  \text{angles per matrix} = \frac{D}{2}\log_2 D = 32 \times 6 = 192.
\end{equation}
This is trained classically via PyTorch autograd and provides a structured
approximation to the full $O(64)$ group (analogous to an FFT butterfly).
\textbf{Trainable parameters.}
$T_{\mathrm{Givens}} = 192 \times 2 \times 4 = 1{,}536$.

The Givens butterfly cannot be directly deployed on quantum hardware
(it is a classical differentiable parameterization of an orthogonal matrix,
not a gate-model circuit), but provides a strong within-architecture baseline
for comparing the effect of $T$ on generalization.

\subsection{Sham Control}
\label{sec:sham}

The sham replaces both $U$ and $V$ in each attention layer with
unconstrained $\mathrm{Linear}(64, 64, \mathrm{bias=False})$
layers, initialized with random orthogonal matrices but unconstrained
during training.
The sham has \emph{strictly more} parameters than the quantum model
in all configurations (Table~\ref{tab:params}),
so any observed quantum advantage cannot be attributed to capacity.

\begin{table}[h!]
\centering
\caption{Parameter counts for all model variants.}
\label{tab:params}
\begin{tabular}{@{}lrr@{}}
\toprule
Variant & Circuit/Attention Params & Total Params \\
\midrule
RY+CNOT quantum & 384 angles & 142,467 \\
Givens butterfly quantum & 1,536 angles & 143,619 \\
Sham (unconstrained Linear) & $4 \times 2 \times 64^2 = 32{,}768$ & 174,851 \\
Classical MHA (independent Q,K,V) & $12 \times 64^2 = 49{,}152$ & 207,619 \\
\midrule
\textit{MAE-ViT baseline (pretrained)} & \textit{---} & \textit{2,722,947} \\
\bottomrule
\end{tabular}
\end{table}

\subsection{Classical MHA Baseline}

Standard multi-head attention with 8 heads (head dimension $= 8$),
independent learned Q, K, V projections, same transformer depth and MLP structure.
This controls for the effect of the orthogonal constraint at fixed architecture.

% ─────────────────────────────────────────────────────────────
\section{Experiments}
\label{sec:experiments}

\subsection{Experimental Setup}

\paragraph{Datasets.}
We evaluate on three DeepLense simulation datasets~\cite{alexander2020deep},
all $64\times64$ grayscale with three classes (axion, CDM, no substructure):
\begin{itemize}
  \item \textbf{Model\_I}: Euclid-like observation characteristics, weaker axion
        substructure signal (lower SNR); AUC ceiling $\approx 0.96$.
  \item \textbf{Model\_II}: Euclid-like characteristics, stronger axion signal;
        AUC ceiling $\approx 0.99$; \emph{primary evaluation dataset}.
  \item \textbf{Model\_III}: Hubble Space Telescope (HST) observation
        characteristics; AUC ceiling $\approx 0.99$.
\end{itemize}
Each dataset contains ${\approx}25{,}000$ images per class
(${\approx}29{,}500$--$29{,}900$ in the full-data runs).
Train/validation split is 9:1 at \texttt{seed=42}.

\paragraph{Training protocol.}
All experiments use AdamW optimizer
(learning rate $3\times10^{-4}$, weight decay $0.01$,
$(\beta_1,\beta_2) = (0.9,0.999)$) for 50 epochs with cosine annealing.
Batch size is 64.
Loss is cross-entropy on 3-class labels.
The checkpoint with the highest macro one-vs-rest AUC on the validation set
is retained as the final model; metrics are \emph{not} taken from the last epoch.
No data augmentation is applied; images are normalized to $[0,1]$.
Hardware: NVIDIA H200 GPU; implementation: pure PyTorch for both Givens and
RY+CNOT (see Section~\ref{sec:rycnot}), without PennyLane overhead.

\paragraph{Evaluation metric.}
Primary metric: macro-averaged one-vs-rest AUC (each of the three classes is
treated as the positive class against the other two; AUC is averaged).
This follows the DeepLense benchmark convention~\cite{alexander2020deep,prasha2025masked}.

\paragraph{Notation.} $\dagger$ marks cells with $\geq2$ independent runs at
identical configuration, reported as mean (range); this makes run-to-run
variance visible rather than reporting a single cherry-picked value.

\subsection{Full-Data Results}
\label{sec:fulldata}

\begin{table}[h!]
\centering
\caption{Full-data AUC (${\approx}25{,}000$/class, 9:1 split, \texttt{seed=42}).
MAE-ViT is our own reproduction under the protocol of~\cite{prasha2025masked}.
$\dagger$: mean (range) over multiple single-seed runs at identical config.
Best per dataset in \textbf{\textcolor{good}{green}}.}
\label{tab:fulldata}
\begin{tabular}{@{}lccccc@{}}
\toprule
Method & Params & Model\_I & Model\_II & Model\_III \\
\midrule
RY+CNOT quantum   & 142,467 & 0.979$^\dagger$ & 0.985$^\dagger$ (0.976--0.997) & 0.995$^\dagger$ \\
RY+CNOT sham      & 174,851 & 0.9801 & 0.991$^\dagger$ & 0.9903 \\
Givens quantum    & 143,619 & \win{0.9813} & \win{0.9940} & \win{0.9962} \\
Givens sham       & 174,851 & 0.9803 & 0.9910 & 0.9887 \\
Classical MHA     & 207,619 & 0.9807 & 0.9921 & 0.9957 \\
MAE-ViT (ours)~\cite{prasha2025masked} & 2,722,947 & 0.9778 & 0.9895 & 0.9895 \\
\bottomrule
\end{tabular}
\end{table}

At full data, the four architectures (RY+CNOT, Givens, sham, classical MHA)
cluster tightly within $\approx$1 percentage point of each other on Models
I and III; on Model\_II, RY+CNOT's mean (0.985) is nominally below classical
MHA (0.9921) and both sham controls, though its 3-run range (0.976--0.997)
overlaps them.
This is a meaningfully different picture from the sham-vs-quantum framing
in Section~\ref{sec:sweep}: at full data, no architecture in this comparison
shows an advantage that survives its own run-to-run variance.
All four architectures, plus our reproduced MAE-ViT (0.978--0.990), comfortably
exceed 0.97 AUC at full data---the interesting regime for distinguishing
these methods is the low-$N$ sweep (Section~\ref{sec:sweep}), not the
full-data ceiling.

\subsection{Data-Efficiency Sweep: Primary Result}
\label{sec:sweep}

\begin{table}[h!]
\centering
\caption{AUC vs.\ $N$/class on Model\_II (\texttt{seed=42}).
$^\dagger$: mean over multiple runs, individual values in parentheses.
$^\ddagger$: run failed to train (see text); value is the failed run's own
best-checkpoint AUC, included for completeness rather than excluded.
Highlighted row: the pivotal $N{=}500$ comparison.}
\label{tab:sweep}
\begin{tabular}{@{}rcccc@{}}
\toprule
$N$/class & RY+CNOT quantum & RY+CNOT sham & Givens quantum & Classical MHA \\
\midrule
\rowcolor{rowhi}
500  & 0.9144 & 0.8812 & 0.526$^\dagger$ (0.513, 0.539) & 0.910$^\dagger$ (0.908, 0.912) \\
3000 & 0.9638 & 0.9595 & 0.9389 & 0.964$^\dagger$ (0.962, 0.966) \\
5000 & 0.9700 & 0.9613 & 0.9613 & 0.968$^\dagger$ (0.968, 0.968) \\
8000 & 0.9662 & 0.9653 & 0.558$^\ddagger$ & 0.980$^\dagger$ (0.978, 0.983) \\
full & 0.985$^\dagger$ (0.976, 0.983, 0.997) & 0.991$^\dagger$ & \win{0.9940} & 0.9921 \\
\bottomrule
\end{tabular}
\end{table}

The headline $N{=}500$ comparison must be read against \emph{two} baselines,
not one.
RY+CNOT (AUC $0.9144$) does clearly beat Givens (AUC $0.526$ mean, near
chance)---but it is statistically indistinguishable from classical MHA
(AUC $0.910$ mean, $\Delta = +0.004$, well within single-seed noise).
This changes the paper's central claim: the $N{=}500$ gap is not evidence
that \emph{quantum parameterization beats classical attention}: it is evidence
that \emph{the Givens butterfly specifically fails to train reliably at small
$N$}, while RY+CNOT, its sham control, and unconstrained classical MHA all
train successfully and land within $\approx$1 point of each other.

\subsubsection{Is the Givens Failure Memorization or Optimization Failure?}
\label{sec:caro-prediction}

The original hypothesis---that Givens' $T{=}1{,}536$ angles memorize the
$N{=}500{\times}3{=}1{,}500$-example training set---predicts near-perfect
\emph{training} accuracy alongside near-chance validation AUC.
We did not log training accuracy for the reported runs, so we cannot confirm
this directly; however, per-epoch validation curves for the failed runs
(e.g.\ Model\_II, $N{=}8000$, Section~\ref{sec:crossover}) show
validation AUC oscillating near $0.50$--$0.55$ for all 50 epochs with a final
confusion matrix that predicts a single class for every input---behavior
consistent with the optimizer never escaping a degenerate region of the loss
landscape, not with a model that fit the training set and failed to
generalize.
We regard the "memorization" account as unconfirmed and the "optimization
instability in the butterfly parameterization" account as the more
parsimonious explanation given the available evidence; distinguishing them
rigorously requires logging train-set accuracy, which we leave for future
work (Section~\ref{sec:limitations}).

For completeness, the Caro bound values at $N{=}500$ are:
\begin{equation}
  \epsilon_{\mathrm{Givens}}
  \;\leq\; C\sqrt{\frac{1536}{500}} \approx 1.75\,C,
  \quad
  \epsilon_{\mathrm{RY}}
  \;\leq\; C\sqrt{\frac{384}{500}}  \approx 0.88\,C.
  \label{eq:bound-500}
\end{equation}
This bound predicts a \emph{softer} degradation for Givens (larger worst-case
generalization gap), not the catastrophic, unstable failures we observe;
a bound violation of this magnitude is more consistent with an optimization
pathology outside the bound's assumptions (e.g.\ non-convex landscape
with poor conditioning) than with the bound's own mechanism.
We caution against over-crediting the Caro bound for a phenomenon that may
be architecture-specific optimization behavior.

\subsection{Instability, Not Recovery}
\label{sec:crossover}

We initially described Givens' behaviour as $N$ increases as "recovery":
$N{=}3000 \to 0.9389$; $N{=}5000 \to 0.9613$; $N{=}\text{full} \to 0.9940$.
This monotone-looking sequence is misleading in isolation: at $N{=}8000$,
between the ``recovered'' $N=5000$ and full-data points, Givens
\emph{collapses again} to AUC $0.558$ (Table~\ref{tab:sweep}), with the same
single-class-prediction failure mode as the $N{=}500$ runs
(Section~\ref{sec:caro-prediction}).
The correct description is not a smooth crossover governed by a
generalization-expressivity trade-off, but \emph{intermittent training
instability that can occur at any $N$}, whose probability appears to
decrease (not vanish) as $N$ grows.
Of the 6 Givens quantum runs we logged across $N \in \{500, 3000, 5000,
8000, \text{full}\}$ (with 2 seeds at $N{=}500$), 3 failed
(both $N{=}500$ seeds, and $N{=}8000$)---a $50\%$ failure rate in this
sample.
RY+CNOT and classical MHA show no comparable failures in any logged run at
any $N$.
We do not have a confirmed mechanism for why the Givens butterfly
parameterization is more failure-prone; a plausible hypothesis is that its
deeper compositional structure ($\log_2 D = 6$ sequential Givens rotation
stages per matrix) creates a more ill-conditioned optimization landscape
than RY+CNOT's shallower, more redundant ring structure, but this is
speculative pending a dedicated loss-landscape study.

\subsection{Circuit Depth Ablation}
\label{sec:depth}

We ablate the number of RY+CNOT layers $L_\mathrm{circ} \in \{4, 8, 12\}$
on Model\_II to test whether the small-$N$ advantage is governed by the
Caro bound, which (via Corollary~\ref{cor:design}, $T$ linear in
$L_\mathrm{circ}$) predicts $L{=}4$ ($T{=}192$) should \emph{outperform}
$L{=}8$ ($T{=}384$) at $N{=}500$, since a smaller $T$ gives a tighter bound.

\begin{table}[h!]
\centering
\caption{Circuit depth ablation on Model\_II, RY+CNOT quantum mode
(single seed per cell).}
\label{tab:depth}
\begin{tabular}{@{}rrrcc@{}}
\toprule
$L_\mathrm{circ}$ & Angles/matrix & Total $T$ & AUC @ $N{=}500$ & AUC @ full \\
\midrule
4  & 24 & 192 & 0.8961 & 0.9822 \\
\rowcolor{rowhi}
8  & 48 & 384 & \win{0.9144} & 0.985$^\dagger$ \\
12 & 72 & 576 & \win{0.9298} & 0.9810 \\
\bottomrule
\end{tabular}
\end{table}

\textbf{The Caro-bound prediction is refuted by this data.}
AUC increases \emph{monotonically} with $L_\mathrm{circ}$ at $N{=}500$
($0.8961 \to 0.9144 \to 0.9298$), the opposite of what a bound in which
smaller $T$ implies better small-$N$ generalization would predict.
At full data, the three depths are statistically indistinguishable
($0.981$--$0.985$).
The most direct reading is that within the RY+CNOT family, more circuit
layers give the optimizer more expressive power to fit useful angular
features \emph{without} the instability that afflicts the much larger,
differently-structured Givens parameterization---i.e., for RY+CNOT,
capacity is still expressivity-limited rather than generalization-limited
even at $N{=}500$, and the Caro bound's small-$T$-is-better logic does not
bind in this range of $T \in [192, 576]$.
This is consistent with our conclusion in Section~\ref{sec:crossover} that
the $N{=}500$ Givens failure is better explained by optimization instability
specific to the butterfly parameterization than by a generalization-bound
effect that should also, and did not, appear in the RY+CNOT depth sweep.

\subsection{Training Audit}
\label{sec:audit}

We verify circuit contribution and trainability for RY+CNOT at $N{=}500$:

\begin{table}[h!]
\centering
\caption{Circuit health diagnostics (RY+CNOT quantum, Model\_II, $N{=}500$).}
\label{tab:audit}
\begin{tabular}{@{}lcc@{}}
\toprule
Diagnostic & Value & Interpretation \\
\midrule
Gradient norm at circuit params (epoch 20--50) & $> 0$ (stable) & No barren plateau \\
Val-AUC trajectory & monotone to $0.9144$ & Circuit is learning \\
Expressibility ($\hat{\Delta}_{KL}$ vs.\ Haar) & moderate & Not over-entangled \\
\bottomrule
\end{tabular}
\end{table}

The absence of barren plateaus in the $L{=}8$, $N_Q{=}6$ configuration
is consistent with prior results on structured, shallow circuits~\cite{cerezo2021cost}:
nearest-neighbor entanglement and moderate depth keep gradient norms
measurably nonzero throughout training.

\subsection{Comparison with Companion Paper QVF-Scratch}
\label{sec:comparison}

\begin{table}[h!]
\centering
\caption{QOVT variants vs.\ QVF-Scratch on Model\_II.
$^\dagger$: mean over multiple runs (see Table~\ref{tab:sweep}/\ref{tab:fulldata}
for individual values and ranges).}
\label{tab:comparison}
\begin{tabular}{@{}lrrrcc@{}}
\toprule
Method & $T$ & Total Params & Full-data AUC & $N{=}500$ AUC & NISQ \\
\midrule
QVF-Scratch quantum~\cite{wang2025qvf} & 96 & 142,795 & \win{0.9983} & ${\approx}0.97$ & No \\
QOVT RY+CNOT quantum & 384 & 142,467 & 0.985$^\dagger$ & 0.9144 & \textbf{Yes} \\
QOVT Givens quantum & 1,536 & 143,619 & 0.9940 & 0.526$^\dagger$ (unstable) & No \\
Classical MHA & --- & 207,619 & 0.9921 & 0.910$^\dagger$ & No \\
MAE-ViT (ours)~\cite{prasha2025masked} & --- & 2,722,947 & 0.9895 & --- & No \\
\bottomrule
\end{tabular}
\end{table}

QVF-Scratch (companion paper~\cite{wang2025qvf}) achieves the highest
full-data AUC ($0.9983$) via amplitude encoding, but uses a CNN encoder
and is simulation-only.
QOVT RY+CNOT is comparable in full-data AUC to classical MHA and Givens
(all within $\approx$1 point, Table~\ref{tab:fulldata}) but is
\textbf{NISQ-deployable}, embedding the quantum circuit directly within
the transformer attention mechanism at a depth ($O(48)$) compatible with
current hardware coherence times.
Unlike QVF-Scratch, whose 21/22-positive-margin sweep provides evidence
for a consistent quantum-over-sham advantage, QOVT's own sweep
(Table~\ref{tab:sweep}) shows RY+CNOT statistically tied with classical MHA
at every tested $N$; QOVT's primary contribution is therefore
\textbf{hardware deployability at competitive accuracy}, not a demonstrated
quantum advantage over classical attention.

% ─────────────────────────────────────────────────────────────
\section{Analysis}
\label{sec:analysis}

\subsection{What the Orthogonality Constraint Does and Does Not Explain}

The orthogonality constraint $U^\top U = I$ provides two properties we can
state with confidence, and one claim that our own data does not support.

\paragraph{1.\ Parameter efficiency (established).}
$U$ and $V$ are each specified by $T/4$ rotation angles rather than $D^2$
free entries---a real, architecture-level reduction, independent of whether
it translates into a generalization advantage.

\paragraph{2.\ Spectral stability (established).}
All singular values of $U$ and $V$ are identically $1$, so
$\|U\|_2 = \|V\|_2 = 1$.
Gradient magnitudes through the orthogonal projections are exactly preserved,
which plausibly explains why RY+CNOT never exhibited the catastrophic
training failures we observed for Givens (Section~\ref{sec:crossover}):
its angles cannot individually blow up a key vector's norm the way an
unconstrained linear layer's singular values can.

\paragraph{3.\ Caro-bound tightening (not supported at the tested $T$ values).}
We had hypothesized that $T{=}384$'s tighter generalization bound relative to
classical MHA's unconstrained parameterization would produce a measurable
small-$N$ advantage.
The data do not support this: RY+CNOT and classical MHA are statistically
tied at every tested $N$ (Table~\ref{tab:sweep}), and the depth ablation
(Table~\ref{tab:depth}) shows AUC \emph{increasing} with $T$ at $N{=}500$
within the RY+CNOT family---the opposite of what a binding Caro bound would
predict.
We do not have evidence that the orthogonality constraint's contribution to
this benchmark operates primarily through the generalization-bound channel;
the spectral-stability channel (item 2) is the more empirically supported
mechanism given that it correctly predicts \emph{which} circuit family
(RY+CNOT, not Givens) avoids training failure.

\subsection{Why Givens Fails Intermittently}

We initially hypothesized that Givens' $T{=}1{,}536$ angles memorize the
small $N{=}500$ training set, producing near-chance generalization.
This account cannot explain why Givens \emph{also} fails at $N{=}8000$
(Section~\ref{sec:crossover}), where $24{,}000$ total training examples are
far too many to memorize with $1{,}536$ parameters.
A memorization account predicts failure probability decreasing monotonically
and vanishing well before $N{=}8000$; we instead observe failure at both
the smallest and one of the largest tested $N$, with no failures at
intermediate $N \in \{3000, 5000\}$ in our (small) sample.
This pattern is more consistent with \textbf{stochastic optimization
failure}---the training run occasionally lands in a degenerate region of a
poorly-conditioned loss landscape and never escapes it within 50 epochs,
independent of $N$---than with a data-size-dependent generalization
phenomenon.
The deeper compositional structure of the butterfly ($\log_2 D = 6$
sequential Givens-rotation stages composed per matrix, versus RY+CNOT's
shallower, redundant ring) is a plausible source of this conditioning
sensitivity, but we have not directly measured the loss landscape to confirm
it.

\subsection{Revisiting the Full-Data Comparison}

At full data, we no longer find clear separation between architectures:
RY+CNOT (mean $0.985$, range $0.976$--$0.997$ across 3 runs), Givens
($0.9940$, 1 run), sham controls ($0.987$--$0.991$), classical MHA
($0.9921$), and our reproduced MAE-ViT ($0.9895$) are within
$\approx1.5$ points of each other on Model\_II, and RY+CNOT's own
run-to-run range alone spans $2$ points.
We no longer describe this as a "crossover" governed by an
expressivity-vs-generalization trade-off (our earlier framing); the data are
better summarized as: \emph{multiple architecturally distinct approaches
converge to similar full-data performance on this benchmark}, and the
differences that do appear at small $N$ are driven primarily by
Givens-specific training instability rather than a general
principle about quantum circuit design.

% ─────────────────────────────────────────────────────────────
\section{Conclusion}
\label{sec:conclusion}

We applied Quantum Orthogonal Vision Transformers (QOVT) to three-class dark
matter substructure classification on the DeepLense gravitational lensing
benchmark, comparing two quantum circuit families---RY+CNOT ring ($T{=}384$)
and Givens butterfly ($T{=}1{,}536$)---against classical sham and
multi-head-attention controls, with all runs now complete
(no \running{} entries remain).

\paragraph{Summary of findings.}
The headline empirical result---RY+CNOT beating Givens by $0.40$ AUC at
$N{=}500$---survives scrutiny, but the interpretation we originally attached
to it does not.
Classical MHA (unconstrained, more parameters) matches RY+CNOT within
noise at every tested $N$ from 500 to full data (Table~\ref{tab:sweep}),
so the $N{=}500$ gap is evidence of \emph{Givens-specific training
instability}, not of quantum circuits outperforming classical attention.
Corroborating this: (i) Givens fails again at $N{=}8{,}000$
(AUC $0.558$), a data size too large for its parameter count to memorize,
contradicting the memorization account we initially proposed;
(ii) the RY+CNOT depth ablation shows AUC \emph{increasing} monotonically
with $T \in \{192, 384, 576\}$ at $N{=}500$---the opposite of the Caro-bound
prediction that motivated the ablation.
At full data, RY+CNOT, Givens, both sham controls, classical MHA, and our
reproduced MAE-ViT baseline (AUC $0.9895$, corrected from an earlier,
erroneously low reported value) all land within $\approx1.5$ points of
each other on Model\_II.

We regard the confirmed contribution of this paper as narrower than
originally framed: RY+CNOT is a \textbf{NISQ-deployable circuit that matches
classical attention's accuracy on this benchmark at every tested data size},
which is a meaningful engineering result (a working quantum circuit that
does not underperform its classical counterpart) but is not evidence of a
quantum learning advantage governed by the Caro bound.
The Givens butterfly's intermittent, severe training failures (50\% of
logged runs at $N \in \{500, 8000\}$) are a genuine and unexplained
phenomenon worth reporting, but the evidence points to optimization
instability in its deeper compositional structure rather than a
generalization-theoretic effect.

\paragraph{Hardware deployability.}
The RY+CNOT ring circuit (8 layers of RY + CNOT ring on 6 qubits) is
natively executable on IBM Quantum, IonQ, and other superconducting/trapped-ion
architectures with ring or linear-chain connectivity.
Its circuit depth $O(48)$ is well within current coherence windows,
and all trained parameters transfer directly from simulation to hardware.
This remains, to our knowledge, a working quantum attention architecture for
astrophysical image classification that is ready for real hardware execution,
independent of whether it demonstrates a quantum \emph{advantage}.

\paragraph{Limitations.}
\label{sec:limitations}
All results are single-seed except where explicitly noted as multi-run in
Tables~\ref{tab:fulldata}--\ref{tab:comparison}; even the multi-run cells use
only 2--3 seeds, insufficient for tight confidence intervals.
We did not log training-set accuracy, so the memorization-vs-optimization-failure
question for Givens (Section~\ref{sec:caro-prediction}) remains unresolved.
The asymmetric $W_Q{=}I$, $W_K{=}U$ attention design (Section~\ref{sec:qoattn})
differs from the symmetric $W_Q{=}W_K{=}U$ design our original theoretical
framing assumed; we have not re-run experiments with the symmetric design,
so we cannot rule out that it behaves differently.
Real hardware execution with noise characterization has not yet been
performed.

\paragraph{Future work.}
Priority directions, revised in light of the above:
(1)~multi-seed (5+) reruns of the Givens $N{=}500$ and $N{=}8{,}000$
    conditions with train-accuracy logging, to directly test the
    memorization vs.\ optimization-failure hypotheses;
(2)~implement and test the symmetric $W_Q{=}W_K{=}U$ attention variant that
    the original theoretical framing described, to check whether it
    behaves differently from the asymmetric variant we report here;
(3)~a loss-landscape or Hessian-conditioning study comparing Givens and
    RY+CNOT to test the optimization-instability hypothesis directly;
(4)~real-hardware deployment on IBM Quantum Eagle/Heron with zero-noise
    extrapolation~\cite{bharti2022nisq}, now that the RY+CNOT accuracy
    claim (matching, not exceeding, classical MHA) is on firmer footing.

\paragraph{Broader Impact.}
This work is a cautionary as well as a positive result: our first-pass
analysis over-attributed an accuracy gap to quantum generalization theory
before checking it against a classical baseline (MHA) and a depth ablation
that turned out to refute the theoretical prediction.
Both checks were cheap and available before submission; we include this
correction as a concrete illustration of why quantum-advantage claims in
NISQ-era ML need matched classical baselines and falsifiable, pre-registered
predictions, not just favorable comparisons against one classical control.
No harmful applications are anticipated.

% ─────────────────────────────────────────────────────────────

\begin{thebibliography}{99}

\bibitem{alexander2020deep}
S.~Alexander et al.
\textit{Deep Learning the Morphology of Dark Matter Substructure}.
The Astrophysical Journal, 893(1):15, 2020.

\bibitem{brehmer2019mining}
J.~Brehmer, S.~Mishra-Sharma, J.~Hermans, G.~Louppe, and K.~Cranmer.
\textit{Mining for Dark Matter Substructure: Inferring Subhalo Population
Properties from Strong Lenses with Machine Learning}.
The Astrophysical Journal, 886(1):49, 2019.

\bibitem{caro2022}
M.~C.~Caro et al.
\textit{Generalization in Quantum Machine Learning from Few Training Data}.
Nature Communications, 13:4919, 2022.

\bibitem{cerezo2021cost}
M.~Cerezo, A.~Sone, T.~Volkoff, L.~Cincio, and P.~J.~Coles.
\textit{Cost Function Dependent Barren Plateaus in Shallow Parametrized
Quantum Circuits}.
Nature Communications, 12:1791, 2021.

\bibitem{cherrat2022}
E.~A.~Cherrat, I.~Kerenidis, N.~Mathur, J.~Landman, M.~Strahm,
and Y.~Y.~Li.
\textit{Quantum Vision Transformers}.
Quantum, 8:1265, 2024. arXiv:2209.08167.

\bibitem{comajoan2024quantum}
M.~Comajoan~Cara et al.
\textit{Quantum Vision Transformers for Quark--Gluon Classification}.
MDPI Axioms, 13(5):323, 2024.

\bibitem{dalal2002detecting}
N.~Dalal and C.~S.~Kochanek.
\textit{Direct Detection of Cold Dark Matter Substructure}.
The Astrophysical Journal, 572:25--33, 2002. arXiv:astro-ph/0111456.

\bibitem{dosovitskiy2021image}
A.~Dosovitskiy et al.
\textit{An Image is Worth 16x16 Words: Transformers for Image Recognition
at Scale}.
ICLR 2021. arXiv:2010.11929.

\bibitem{fei2022ovit}
Y.~Fei, Y.~Liu, X.~Wei, and M.~Chen.
\textit{O-ViT: Orthogonal Vision Transformer}.
arXiv:2201.12133, 2022.

\bibitem{gao2018quantum}
X.~Gao, Z.-Y.~Zhang, and L.-M.~Duan.
\textit{A Quantum Machine Learning Algorithm Based on Generative Models}.
Science Advances, 4(12):eaat9004, 2018.

\bibitem{havlicek2019supervised}
V.~Havl\'{i}\v{c}ek et al.
\textit{Supervised Learning with Quantum-Enhanced Feature Spaces}.
Nature, 567(7747):209--212, 2019.

\bibitem{hezaveh2017fast}
Y.~D.~Hezaveh, L.~P.~Levasseur, and P.~J.~Marshall.
\textit{Fast Automated Analysis of Strong Gravitational Lenses with
Convolutional Neural Networks}.
Nature, 548:555--557, 2017.

\bibitem{hui2017ultralight}
L.~Hui, J.~P.~Ostriker, S.~Tremaine, and E.~Witten.
\textit{Ultralight Axion Dark Matter}.
Physical Review D, 95(4):043541, 2017. arXiv:1610.08297.

\bibitem{lanusse2018cmu}
F.~Lanusse et al.
\textit{CMU DeepLens: Deep Learning for Automatic Image-Based Galaxy--Galaxy
Strong Lens Finding}.
Monthly Notices of the Royal Astronomical Society, 473(3):3895--3906, 2018.

\bibitem{huang2018orthogonal}
L.~Huang, X.~Liu, B.~Lang, A.~W.~Yu, Y.~Wang, and B.~Li.
\textit{Orthogonal Weight Normalization: Solution to Optimization over Multiple
Dependent Stiefel Manifolds in Deep Neural Networks}.
AAAI, 2018. arXiv:1709.06079.

\bibitem{huang2021power}
H.-Y.~Huang et al.
\textit{Power of Data in Quantum Machine Learning}.
Nature Communications, 12:2631, 2021.

\bibitem{kerenidis2021classical}
I.~Kerenidis, J.~Landman, and N.~Mathur.
\textit{Classical and Quantum Algorithms for Orthogonal Neural Networks}.
arXiv:2106.07198, 2021.

\bibitem{landman2022quantum}
J.~Landman et al.
\textit{Quantum Methods for Neural Networks and Application to Medical Image
Classification}.
Quantum, 6:881, 2022. arXiv:2212.07389.

\bibitem{li2022quantum}
G.~Li, X.~Zhao, and X.~Wang.
\textit{Quantum Self-Attention Neural Networks for Text Classification}.
arXiv:2205.05625, 2022.

\bibitem{liu2021rigorous}
Y.~Liu, S.~Arunachalam, and K.~Temme.
\textit{A Rigorous and Robust Quantum Speed-Up in Supervised Machine Learning}.
Nature Physics, 17(9):1013--1017, 2021.

\bibitem{mcclean2018barren}
J.~R.~McClean, S.~Boixo, V.~N.~Smelyanskiy, R.~Babbush, and H.~Neven.
\textit{Barren Plateaus in Quantum Neural Network Training Landscapes}.
Nature Communications, 9:4812, 2018.

\bibitem{ngairangbam2022anomaly}
V.~S.~Ngairangbam, M.~Spannowsky, and M.~Takeuchi.
\textit{Anomaly Detection in High-Energy Physics Using a Quantum Autoencoder}.
Physical Review D, 105:095004, 2022.

\bibitem{pennylane}
V.~Bergholm et al.
\textit{PennyLane: Automatic Differentiation of Hybrid Quantum-Classical
Computations}.
arXiv:1811.04968, 2018.

\bibitem{prasha2025masked}
A.~A.~Prasha et al.
\textit{Masked Autoencoder Pretraining on Strong-Lensing Images for Joint
Dark-Matter Model Classification and Super-Resolution}.
arXiv:2512.06642, 2025.

\bibitem{preskill2018quantum}
J.~Preskill.
\textit{Quantum Computing in the NISQ Era and Beyond}.
Quantum, 2:79, 2018.

\bibitem{bharti2022nisq}
K.~Bharti et al.
\textit{Noisy intermediate-scale quantum ({NISQ}) algorithms}.
Reviews of Modern Physics, 94(1):015004, 2022.

\bibitem{sim2019expressibility}
S.~Sim, P.~D.~Johnson, and A.~Aspuru-Guzik.
\textit{Expressibility and Entangling Capability of Parameterized Quantum
Circuits for Hybrid Quantum-Classical Algorithms}.
Advanced Quantum Technologies, 2(12):1900070, 2019.

\bibitem{terashi2021event}
K.~Terashi et al.
\textit{Event Classification with Quantum Machine Learning in
High-Energy Physics}.
Computing and Software for Big Science, 5:2, 2021.

\bibitem{tesi2024quantum}
A.~Tesi et al.
\textit{Quantum Attention for Vision Transformers in High Energy Physics}.
arXiv:2411.13520, 2024.

\bibitem{unlu2024hybrid}
E.~B.~Unlu et al.
\textit{Hybrid Quantum Vision Transformers for Event Classification in
High Energy Physics}.
MDPI Axioms, 13(3):187, 2024.

\bibitem{wang2025qvf}
S.~Wang, C.~Theobalt, and V.~Golyanik.
\textit{Quantum Visual Fields with Neural Amplitude Encoding}.
NeurIPS 2025. arXiv:2508.10900.

\end{thebibliography}

\end{document}
